\documentclass[11pt]{article}
\usepackage[a4paper,margin=1in]{geometry}
\usepackage{parskip}
\usepackage[hidelinks]{hyperref}
\pagestyle{empty}
\setlength{\parskip}{0.7em}

\begin{document}

\noindent Kurniabudi \emph{(Corresponding Author)}\\
Master of Information System, Universitas Dinamika Bangsa\\
Jambi 36138, Indonesia\\
\href{mailto:kurniabudi@unama.ac.id}{kurniabudi@unama.ac.id}\\
\emph{Manuscript submitted by co-author Eko Arip Winanto on behalf of all authors.}

\vspace{1em}
\noindent To: The Editor-in-Chief\\
\noindent \emph{IEEE Transactions on Information Forensics and Security}

\vspace{1em}
\noindent \textbf{Re: Submission of an original research manuscript}

\vspace{0.5em}
\noindent Dear Editor,

We are pleased to submit our manuscript, \textbf{``Surviving Sanitization with
QuantGhost: Deployment-Robust and Stealthy Steganography in Neural-Network Weights,''}
for consideration as an original research article in the \emph{IEEE Transactions on
Information Forensics and Security}.

Pre-trained neural networks are increasingly shared through public model hubs, which
turns the model file itself into a supply-chain attack surface. Prior weight-steganography
attacks (StegoNet, EvilModel, MaleficNet) hide megabyte-scale payloads in float32
weights, but, as we demonstrate on real models, those payloads are completely erased by
the routine deployment transforms that production systems apply: INT8 quantization,
pruning, and fine-tuning. These are also the very operations the community has proposed
as defenses. Our paper asks, and answers, whether a payload can survive the full
deployment pipeline while remaining stealthy and accuracy-preserving, and turns the answer
into both a fundamental limit and a concrete defense.

The manuscript makes the following contributions:

\begin{itemize}
\item \textbf{A capacity--robustness--stealth trilemma.} We prove that no
  weight-embedding scheme can be simultaneously high-capacity, permutation-invariant,
  and stealthy-quantization-robust, with a tight bound of $\sum_\ell k_h^{(\ell)}
  k_w^{(\ell)}$ bits on any permutation-invariant channel. This limit governs prior
  robust watermarking as well as our own attack.
\item \textbf{QuantGhost.} We present a quantization-aware, distribution-preserving
  embedding that reaches the high-capacity, stealthy, quantization-robust corner of the
  trilemma. It survives INT8, pruning, and fine-tuning at $3.2$\,MB on a real ResNet-50
  for only $0.19$ percentage points of accuracy, and we validate that it survives
  byte-exactly through PyTorch's \emph{production} INT8 kernel, not merely an analytical
  quantizer.
\item \textbf{An attack and its detector.} Against the released Model X-Ray steganalyzer
  QuantGhost is statistically indistinguishable from benign weights on a diverse model
  population, where the prior attack is detected. We then propose a targeted sub-quantum
  residual detector that reliably flags the basic attack across a narrow and a diverse
  $3000$-model zoo, while honestly showing that an adaptive attacker reaches chance on a
  realistic population, so detection must be backed by a structural bound.
\item \textbf{Limits of permutation defenses.} A faithful, full-model neuron permutation
  on a real ResNet-18, verified function-preserving to machine precision over $30$
  independent permutations, removes the high-capacity payload, yet a $196$-bit
  permutation-invariant channel survives and end-to-end bootstraps full recovery of an
  AES-encrypted payload, so permutation-based sanitization is incomplete.
\item \textbf{Generality and reproducibility.} We show the attack is architecture- and
  domain-agnostic on a real Vision Transformer and through a cross-dataset study on
  IoMT/IoT intrusion-detection models (CICIoMT, CICIoT), set out the requirements for a
  complete defense, and release the full code, data, and a one-to-one table/figure-to-script
  mapping on Zenodo.
\end{itemize}

We believe the work is a strong fit for the journal. \emph{IEEE Transactions on
Information Forensics and Security} is a leading venue for steganography, steganalysis,
information hiding, and machine-learning security; our manuscript contributes a proven
information-theoretic limit, a new attack, and a new detector for the protection of
machine-learning supply chains, with direct relevance to the IoT/IoMT deployments TIFS
readers study.

This research is offensive in nature but carried out to strengthen defenses. All
payloads in our experiments are encrypted random bytes or an AES key; no live malware is
embedded, executed, or distributed. We will responsibly disclose our findings to the
authors of the affected detector and defense (Model X-Ray and NeuPerm) prior to
publication, and the paper proposes concrete mitigations.

We confirm that this manuscript is original, has not been published previously, and is
not under consideration for publication elsewhere. All authors have read and approved
the submission and declare no competing interests.

Thank you for considering our work. We look forward to the reviewers' comments.

\vspace{1em}
\noindent Sincerely,

\vspace{1.5em}
\noindent Kurniabudi\\
\emph{Corresponding author, on behalf of all authors}\\
(Eko Arip Winanto, Mohd Yazid Idris, and Kurniabudi)\\[0.3em]
\emph{Submitted via ScholarOne Manuscripts by Eko Arip Winanto
(\href{mailto:winanto@graduate.utm.my}{winanto@graduate.utm.my}).}

\end{document}
